\section{Fine tuning and an explicit instability}\label{sec:instability}
A small perturbation being larger than a closing gap does not by
itself prove instability. Here a concrete perturbation has a
quantitative lower bound in every sector. For each active site's
legal digits $\{0,\ldots,r_j-1\}$ let $P_j$ be the complete-graph
Laplacian, acting as the identity on its side/head symbol. Put
\[
 V_\epsilon=\epsilon\sum_jP_j,\qquad 0<\epsilon\leq1,
\]
and let it vanish off the region. Conditioning on the local classical
radix label makes this one fixed translation-invariant onsite term.
Its per-site norm is at most $\epsilon2^\beta$, so its norm can be
made smaller than any prescribed positive tolerance.

\begin{theorem}[Onsite perturbation destroys the programmed exponent]
\label{thm:instability}
For the two-dimensional bulk ruler and the wetting ruler, this perturbation preserves the
classical selection and the ground energy. In the positive-coupling
phase its thermodynamic gap satisfies, for constants independent of
the selected segment length $\ell$,
\begin{equation}\label{eq:instability}
 \min\left(c_{\rm cl},2\epsilon,
       \frac{2\epsilon\,2^{-\beta}}{5\ell^2}\right)
 \leq\Delta'_\infty(g)\leq\frac{2\pi^2}{\ell^2}.
\end{equation}
Here $c_{\rm cl}$ is the positive classical-sector separation in
the relevant thermodynamic construction. Hence the perturbed gap
still closes, but its coupling exponent is zero.
\end{theorem}
\begin{proof}
The perturbation is positive and annihilates zero-head states with
uniform legal-digit superpositions. Thus every classical block
still has a zero-energy dynamic state, and its classical offset
is unchanged. Legal digits and legal one-head side strings form an
invariant subspace under every counter transition and the new digit
edges. Its complement is entirely penalized and lies at least one
up; exterior errors and multiple heads also cost at least one.

Free zero-head sectors are products of complete graphs, with gap
at least $2\epsilon$ whenever they have a positive eigenvalue.
On the free one-head space, the digit flips connect all legal
digit vectors within each of the $2\ell$ head classes. Write the
counter Laplacian as $L_c$, digit Laplacian as $L_d$, and let $Q$
project to functions constant on each class's digit vectors.
Then $L_d\geq2$ on $Q^\perp$. The compressed quadratic form
$QL_cQ$ contains the cycle
\[
 \INC_0,\ldots,\INC_{\ell-1},
 \RET_{\ell-1},\ldots,\RET_0,\INC_0
\]
with every edge weight at least $2^{-\beta}$. Carry and return
conditions hold on a fraction $1/r_j$ of digit vectors; wrap and
restart close the two ends. Consequently its gap is at least
$\gamma=4\,2^{-\beta}/\ell^2$. Extra increment edges are nonnegative.
This also proves that the enlarged free one-head graph is connected.

For a normalized vector orthogonal to constants write
$\psi=\psi_0+\psi_1$, with $\psi_0=Q\psi$, and let
$x=\|\psi_0\|,y=\|\psi_1\|$. The counter's weighted degree is at
most two, so $\|L_c\|\leq4$. The triangle inequality for
$L_c^{1/2}$ gives
\[
 \langle\psi,(L_c+\epsilon L_d)\psi\rangle
 \geq\max\{2\epsilon y^2,(\sqrt\gamma x-2y)_+^2\}.
\]
If $y\leq\sqrt\gamma x/4$, the second term is at least
$\gamma/8$. Otherwise $y^2\geq\gamma/20$, and the first is
at least $\epsilon\gamma/10$. This proves the one-head lower
bound $2\epsilon2^{-\beta}/(5\ell^2)$. Combining the other
sectors and the classical separation gives the left side of
\eqref{eq:instability}.

For the upper bound take the head-only function
$\cos(\pi(i+1/2)/\ell)$, equal on both rails and constant in
digits. It is orthogonal to constants, digit flips have zero
energy, and the rung edges have zero energy. The remaining path
Rayleigh quotient is at most $2\pi^2/\ell^2$.
Finally $\ell=\Theta(\log(1/g))$ for either ruler. The normalized
logarithm of both bounds tends to zero. Ties do not alter this
estimate when the relevant thermodynamic gap is defined.
\end{proof}

This is an existence theorem for destabilizing perturbations, not a
claim that every perturbation changes the exponent. A proposed
symmetry-protected persistence of the wall geometry under diagonal
wall perturbations requires a new charging estimate at run ends.
It remains a sketch and is not used in any theorem. Making the
corner cost size-dependent to program a structural exponent also
requires a new minimizer and defect analysis; it is not an
established realization of an uncomputable $\nu$.
